Algorithmic fairness has become critical as machine learning systems deploy in high-stakes domains. A common approach enforces fairness constraints during training, such as Demographic Parity (DP)---equal positive prediction rates across protected groups---and Individual Fairness (IF)---similar individuals receiving similar predictions \citep{hardt2016equality,dwork2012fairness}.

Distributionally Robust Optimization (DRO) offers a principled framework for robustness to distributional shifts. DRO-FAIR formulates fairness as a constraint within DRO, learning parameters that satisfy fairness under the worst-case distribution within an uncertainty set \citep{hashimoto2018fairness}.

\textbf{The Gap.} Existing evaluations use \textit{random label noise}---uniformly flipping labels. This does not reflect adversarial scenarios where an attacker with knowledge of the fairness metric strategically corrupts the most impactful samples \citep{solans2021poisoning}. To our knowledge, no prior work computes the \textbf{exact gradient} of fairness metrics with respect to individual labels.

\textbf{Our Contributions.}
(1) \textit{FairnessTargetedPGD}: gradient-based attacks computing
$\partial(\text{fairness})/\partial y_i$ on each label.
(2) The temperature ablation: fixing $\tau=1$ across all $\alpha$ (vs.\
the stepped $\tau{=}100$ schedule for $\alpha\le 0.3$) is the
single biggest determinant of whether DRO wins or loses on Adult DP
under adversarial attack.
(3) The k-NN graph ablation and the empirical-radii interpretation per
Kuldeep's feedback (\S\ref{sec:results-tabular}); IF-attack results are
pending a cluster re-run.
(4) The adversarial-vs-random comparison showing adversarial PGD
raises DP by an order of magnitude more than random label noise at
matched $\alpha$ on Adult.
